\section{形式的対応}
\label{sec:formal}

本節ではVED$\times$IFGTの変数と概念的現象との対応の提案を記録する。
表は形式的役割とモデル内の解釈を分けるものであり、操作的な対応写像を与えない。
Master Variable Table v2.1への帰属には、本論文の冒頭で述べた出典の来歴に関する制限が残る。

\subsection{背景場}
情報密度場$I(x,\tau) = f(1-C(x,\tau))$を、概念構造のモデルにおける連続場として用いる。
本論文の$I$は、新しい認知的原始概念ではなく、IFGTの情報密度という語彙を認知的観測窓で解釈する提案であり、型を明示して確立された同一性ではない。
この局所的な解釈では、$I$は位置$x$、発展パラメータ$\tau$における未解決の情報的緊張の密度を表す。
それは、完全には安定化せず、完全には散逸していない準閉包領域として記述される。

高い$I$には概念的に活動的な領域との対応を、低い$I$には構造が大部分閉包した領域との対応を提案する。
これは較正された活動指標ではない。
本記述では、厳密な正値性$I(x,\tau)>0$を、通常の概念状態としての完全閉包を排するものと解釈する。
それだけで局所的な非ゼロの勾配や流れが保証されるわけではない。
勾配形成における残留非対称性、不均一な履歴、境界構造の役割は、モデルレベルの仮説に留まる。

情報ポテンシャル
\begin{equation}
  \Phi(x,\tau) = (K * I)(x,\tau)
  \label{eq:Phi_primary}
\end{equation}
が主要な幾何的対象である。
これを地形と読むことは、蓄積された偏りに対するモデルレベルの解釈である。
ポテンシャル勾配は有効な方向的偏りを与えるが、流れの関係には拡散項も含まれるため、それだけで完全な軌道を決めない。
場から軌道、軌道から思考への対応は、型を明示した写像や操作的な写像としては与えられていない。

\subsection{変数対応表}
\begin{center}
\renewcommand{\arraystretch}{1.5}
\begin{tabular}{p{0.22\textwidth}p{0.46\textwidth}p{0.22\textwidth}}
\toprule
\textbf{VED$\times$IFGT変数} &
\textbf{概念領域での解釈} &
\textbf{形式的役割} \\
\midrule
$C(x,\tau)$ & $x$における概念的固定の度合い & 背景場；$C \in [0,\Cmax)$ \\
$I(x,\tau) = f(1-C)$ & 未解決の概念的緊張の密度 & 中心となる場；常に$I > 0$ \\
$\Phi = K*I$ & モデル内の概念地形 & ポテンシャル幾何 \\
$\nabla\Phi$ & 認知的抵抗の低い側への偏りの提案 & $\JI$への拘束力 \\
$\JI = -D\nabla I + \mu F_I$ & 思考の表現の提案 & 形式的な流れ場 \\
$\sigma = \sigma_d - \sigma_p - \sigma_b$ & 地形の変形率 & 源項 \\
$\Delta$ & 新たな経験、知覚、言語入力 & 駆動源 \\
$\delta I$ & 未解決の概念的残差 & $I$を構成する痕跡 \\
$\tau$ & 経験的時間との対応の提案 & 形式的な発展・順序パラメータ \\
$\ell_i = \mathcal{Q}[I,\Phi,C;\cdot]$ & 再利用可能な言語トークン（自然言語の単語、記号、紋章、身振り、図、数式、または局所的な摂動として場に作用しうる任意の離散単位） & $\Phi$への作用の提案；シグネチャは未解決 \\
$\Phistar$ & 成熟した沈殿済みの概念地形 & $\Phi$の長時間アトラクタ \\
$\Cmax$ & 理想的な完全閉包境界 & 境界の表現であり、通常の状態ではない \\
\bottomrule
\end{tabular}
\end{center}

時間パラメータは形式的発展を順序づける。その経験的な読みは別の提案であり、時計の較正ではない。
同様に、沈殿した地形と長時間アトラクタという呼称には、第\ref{sec:core}節で論じる極限と履歴依存的解釈の未解決性が残る。

\begin{remark}[記法上の注意]
作用素$\mathcal{Q}$を量子化作用素と呼ぶのは、連続場の配置から離散的で再アドレス可能な単位を抽出するというその作用を、連続変数から離散値への移行と構造的に比較するためである。
これは非形式的な類比であって、数学的な同型や共通の物理的操作ではない。
CT内では「量子」という語を理論的主張として用いないが、この類比は作用素の名称に保持する。
\end{remark}

\subsection{この対応が主張しないこと}
以上の提案的対応は、意味内容、志向性、心理状態を形式体系へ導入するものではない。
特に、
\begin{itemize}
\item $\Phi$は意味場ではない。モデル内の相対的な遷移コストを記述するための動的な拘束幾何である。
\item $\JI$は思考内容ではない。それは形式的な流れ場であり、軌道の記述とその思考としての解釈は異なる使い方である。
経験に現れる内容は下流の主題として扱い、これらの使い方を同一視することで確立される帰結とはしない。
\item $\ell_i$は意味の単位ではない。$\nabla\Phi$へ局所的に結合する、離散的で再アドレス可能な単位である。
\end{itemize}
これらの除外は方法論上の必要事項であり、意味や経験の究極的な本性についての主張ではない。
